\section{Discussion}
\label{sec:discussion}

\noindent\textbf{What the results tell us about the domain gap.}
The 30-pp swing between template-based and topology-preserving
synthetic data is large enough to demand an explanation.
Figure~\ref{fig:graph_stats} provides one: template-generated diagrams
pile up at degree 2--3 because symbols are placed independently and
connected with simple lines, whereas real P\&IDs have a heavy-tailed
degree distribution reflecting the way process systems actually work ---
high-degree junctions at vessel inlets, valve groups on main headers,
instrument clusters around critical equipment.
A model that has never seen these patterns during training will not
learn the right priors for predicting them, regardless of how
realistic the individual symbol images look.
This suggests that the key property of a useful synthetic P\&ID
corpus is not visual fidelity but topological fidelity.

\noindent\textbf{Practical takeaways.}
The 5.6-pp gap between E2 (synth-only) and E1 (oracle) may seem small,
but it is worth pausing on what each configuration represents:
E1 trains on up to 10 real annotated P\&IDs per fold, a resource most
industrial sites cannot legally or practically provide.
E2 trains on zero real P\&IDs and still surpasses the modular baseline
by 18 pp.
From an engineering standpoint, a team that can provide 12 existing
plant drawings as seeds --- without disclosing them externally ---
now has a viable path to a competitive P\&ID digitizer with no additional
annotation effort.
The E3 results show that if even a small number of real images can be
incorporated, performance climbs further, reaching 72.6\% edge mAP and
closing most of the gap to full supervision.

\noindent\textbf{The plateau and what it means.}
The scaling curve in Figure~\ref{fig:scale_ablation} flattens noticeably
beyond around 400 images, with only a 5-pp gain between 400 and 665.
The natural interpretation is that the generator has largely exhausted
the structural variety that 12 seeds can offer.
The same 12 underlying topologies, however many times they are perturbed
and re-rendered, will not introduce new junction configurations or
process flow patterns.
Expanding the seed pool --- say, from 12 to 30 real P\&IDs from
different plant types --- seems likely to shift the plateau substantially
upward and is the most direct path to further improvement.

\noindent\textbf{Limitations.}
Our evaluation rests on 12 images, which introduces variance that
even 15 independent runs cannot fully suppress; results should be
read with the reported standard deviations in mind.
OPEN100 covers a single reactor configuration, so generalisation to
refineries or chemical plants remains unverified.
Finally, obtaining even the seed annotations requires expert time,
though considerably less than annotating a full training corpus.
